\chapter{La solution}
\label{ch:solution}

\section{Le principe : une cascade}

Un seul principe gouverne l'architecture : \textbf{chaque étage ne passe la main au suivant que
s'il ne sait pas faire lui-même}. Comprendre une demande ne demande pas un LLM ; décider d'un
produit ne doit pas en demander un ; mettre en mots un résultat calculé peut en demander un petit ;
lever un vrai doute peut justifier un gros, une fois. C'est la lecture littérale de \og le plus
petit modèle qui suffit \fg, et c'est aussi la réponse à la nuance de CGI sur le \og gros modèle à
un instant précis \fg. Un organisateur l'a résumé ainsi : \og pour gagner, c'est moins de LLM,
plus d'algo \fg. Dans le cas normal, un tour de conversation fait \textbf{au plus un appel à un
LLM, à la fin}.

\begin{figure}[H]
\centering
\begin{adjustbox}{max width=\textwidth}
\begin{tikzpicture}[x=1cm,y=1cm]
  % Acteurs
  \node[acteur,minimum width=3.6cm] (client) at (0,0) {Client, dans l'interface};
  \node[acteur,minimum width=5.4cm,text width=5.2cm] (out) at (0,-9.6)
    {Réponse\\+ panneau \og pourquoi \fg{} + panneau \og frugalité \fg};
  % Étages
  \node[etage,minimum width=4.4cm] (e1) at (0,-1.5)
    {\textbf{1. Comprendre}\\ mots-clés, sinon \f{multilingual-e5-small}\\ + régression logistique ; entités par règles};
  \node[test] (d) at (0,-3.05) {confiance $\geq \tau$ ?};
  \node[etage,minimum width=4.4cm] (e2) at (0,-4.65)
    {\textbf{2. Converser}\\ parcours YAML, profil du client};
  \node[etage,minimum width=3.3cm] (e3) at (-2.55,-6.35)
    {\textbf{3. Conseiller}\\ règles publiques\\ + appétence};
  \node[etage,minimum width=3.3cm] (e4) at (2.55,-6.35)
    {\textbf{4. Agir}\\ outils\\ + confirmation};
  \node[etage,minimum width=4.4cm] (e5) at (0,-8.05)
    {\textbf{5. Formuler}\\ gabarit, sinon \f{ministral-3:3b}\\ + garde-fou};
  % Doute
  \node[module,minimum width=2.9cm] (clar) at (-5.4,-3.05) {Question de\\ clarification};
  \node[risque,minimum width=2.9cm] (esc) at (-5.4,-4.65) {\textbf{Escalade}\\ \f{ministral-3:14b}};
  % Données
  \node[donnee,minimum width=2.6cm] (prof) at (-5.6,-6.35) {Profils clients\\ (simulés)};
  \node[donnee,minimum width=2.6cm] (cata) at (-5.6,-7.6) {Catalogue\\ + sources};
  \node[donnee,minimum width=2.6cm] (jour) at (5.6,-6.35) {Journal\\ des actions};
  \node[donnee,minimum width=2.6cm] (glos) at (5.6,-8.05) {Glossaire\\ (fiches vérifiées)};
  % Hors ligne
  \node[module,dashed,minimum width=3.2cm,text width=3.1cm] (off) at (5.95,-1.5)
    {\textit{Hors ligne, une fois}\\ traduction et exemples\\ en français, qui\\ entraînent l'étage 1\\ (gros modèle, DeepEirb)};
  % Mesure
  \node[module,minimum height=9.6cm,minimum width=0.75cm] (mes) at (8.15,-4.8) {};
  \node[rotate=90,font=\scriptsize,align=center] at (mes.center)
    {\textbf{Mesure continue} : latence, tokens, énergie, Go chargés};
  % Flux principal
  \draw[fl] (client) -- (e1);
  \draw[fl] (e1) -- (d);
  \draw[fl] (d) -- node[et,right]{oui} (e2);
  \draw[fl] (d) -- node[et,above]{non} (clar);
  \draw[flp] (clar.north) |- node[et,pos=0.25,left]{question} (client.west);
  \draw[flp] (clar) -- node[et,right]{si toujours ambigu} (esc);
  \draw[fl] (esc) -- (e2);
  \draw[fl] (e2) -- (e3);
  \draw[fl] (e2) -- (e4);
  \draw[fl] (e2) -- node[et,fill=white]{service simple} (e5);
  \draw[fl] (e3) -- (e5);
  \draw[fl] (e4) -- (e5);
  \draw[fl] (e5) -- (out);
  % Liens aux données
  \draw[lien] (prof) -- (e3);
  \draw[lien] (cata) -- (e3.south west);
  \draw[lien] (e4) -- (jour);
  \draw[lien] (e5) -- (glos);
  \draw[flp] (off.west) -- (e1.east);
\end{tikzpicture}
\end{adjustbox}
\caption{La cascade. Les étages en bleu fonctionnent sans gros modèle ; l'escalade (en rouge) est
l'exception, comptée et affichée.}
\label{fig:cascade}
\end{figure}

Le tableau ci-dessous fixe pour chaque étage ce qu'il reçoit, ce qu'il rend et ce qu'il coûte au
plus. Les budgets de latence sont des \emph{cibles} sur un poste de la salle (RTX 3060) ; les
valeurs réelles sont mesurées et reportées dans le pitch.

\begin{center}\small
\begin{tabularx}{\textwidth}{@{}l Y Y l@{}}
\toprule
\textbf{Étage} & \textbf{Reçoit} & \textbf{Rend} & \textbf{Budget cible} \\
\midrule
1. Comprendre & Le message, l'état du parcours en cours & Famille, intention, confiance, entités &
  $<$ 1 ms à $<$ 30 ms \\
2. Converser & La compréhension, le profil & La prochaine question, ou une demande de conseil ou
  d'action & $<$ 5 ms \\
3. Conseiller & Le besoin complété, le profil, le catalogue & Propositions, raisons, produits
  écartés, alertes & $<$ 20 ms \\
4. Agir & Une action confirmée & Un résultat (référence, ticket) & $<$ 20 ms \\
5. Formuler & Les faits structurés & Le texte affiché & 0 ms (gabarit) à $\sim$1,5 s (3B) \\
Escalade & Le message ambigu & Une intention & $\sim$2 à 4 s, rare \\
\bottomrule
\end{tabularx}
\end{center}

\section{Étage 1 --- Comprendre}
\label{sec:nlu}

\paragraph{Premier filtre : les mots-clés.} Avant tout modèle, un lexique d'expressions par famille
(\og perdu ma carte \fg, \og opposition \fg, \og frais \fg, \og virement \fg, \og mettre de côté \fg,
\og appartement \fg, \og prêt étudiant \fg, \og c'est quoi \fg\dots) est cherché dans le message
normalisé (minuscules, accents retirés, formes simples). Coût : moins d'une milliseconde, aucun
modèle chargé. Le lexique part des mots les plus discriminants de chaque famille dans les données
d'entraînement, puis il est relu à la main. Le filtre \textbf{ne conclut que s'il est sûr} : une
seule famille touchée, et soit une phrase courte, soit plusieurs expressions trouvées. Un seul
mot-clé perdu dans une longue phrase ne suffit pas : le message passe au classifieur. Une
expression n'est gardée dans le lexique que si sa précision mesurée sur la validation dépasse
97\,\%. C'est la suggestion de l'organisateur, et elle se mesure : la part des messages que les
mots-clés traitent seuls, et la précision sur cette part.

\paragraph{Classer en deux niveaux.} Le classifieur prédit d'abord une \emph{famille} (14 : les 9
familles Banking77 plus 5 nouvelles), puis l'\emph{intention} précise (85 : les 77 de Banking77
plus 8 nouvelles). La famille suffit à choisir le parcours ; l'intention sert à affiner la réponse.
Ce double niveau absorbe le bruit d'étiquetage de Banking77 : deux intentions voisines mal séparées
appartiennent presque toujours à la même famille.

\paragraph{Le modèle.} Les phrases sont transformées en vecteurs par
\f{intfloat/multilingual-e5-small} (118 millions de paramètres, vecteurs de 384 nombres, environ
0,47 Go en float32), avec le préfixe \f{query:} attendu par la famille E5. Une régression
logistique (scikit-learn) est entraînée sur ces vecteurs : quelques secondes sur CPU, un fichier de
moins d'un mégaoctet. L'encodeur étant multilingue, on l'entraîne sur trois sources à la fois :
Banking77 en anglais, Banking77 traduit en français, et nos exemples français des nouvelles
intentions. Pour expliquer ce qu'il a compris, l'assistant affiche aussi la phrase
d'entraînement la plus proche (\og votre message ressemble à : \emph{j'ai perdu ma carte} \fg),
retrouvée par un produit matrice-vecteur en mémoire : l'idée du RAG, appliquée à la compréhension.

\paragraph{Les nouvelles intentions.}

\begin{center}\small
\begin{tabularx}{\textwidth}{@{}l l Y@{}}
\toprule
\textbf{Intention} & \textbf{Famille} & \textbf{Exemple de message} \\
\midrule
\f{epargne\_projet} & \f{epargne} & \og Je voudrais mettre de côté pour mes enfants \fg \\
\f{credit\_immobilier} & \f{credit} & \og Je veux acheter un appartement \fg \\
\f{credit\_conso} & \f{credit} & \og Il me faudrait un crédit pour une voiture \fg \\
\f{pret\_etudiant} & \f{credit} & \og Comment on fait un prêt étudiant ? \fg \\
\f{faire\_virement} & \f{virement} & \og Je dois envoyer 500\eur{} à mon fils \fg \\
\f{question\_definition} & \f{aide} & \og C'est quoi un TAEG ? \fg \\
\f{insatisfaction} & \f{insatisfaction} & \og Votre appli ne marche jamais ! \fg \\
\f{hors\_sujet} & \f{conversation} & \og Bonjour \fg, \og merci \fg, \og quel temps fait-il ? \fg \\
\bottomrule
\end{tabularx}
\end{center}

Pour chacune, environ 80 exemples d'entraînement sont \emph{générés} une fois, hors ligne, par un
gros modèle, avec des consignes de variété (registre jeune ou âgé, fautes de frappe, style SMS,
phrases longues) ; un membre de l'équipe les relit tous (une vingtaine de minutes). Les 20 exemples
de \emph{test} de chaque intention sont \emph{écrits à la main} par l'équipe : tester un
classifieur sur des phrases produites par le même générateur que son entraînement surestimerait
son score.

\paragraph{Le seuil de confiance $\tau$.} Si la probabilité de la famille prédite est au moins
$\tau$, on continue. Sinon, l'assistant pose une question fermée construite à partir des deux
familles les plus probables : \og Vous parlez de votre carte, ou d'un virement ? \fg. Si la
réponse reste ambiguë, et seulement alors, le message part au modèle d'escalade. $\tau$ est choisi
par balayage sur un jeu de validation, comme dans Timso AI : on retient la plus petite valeur pour
laquelle la précision sur les messages acceptés dépasse 97\,\%. Une erreur de parcours coûte plus
qu'une question.

\paragraph{Le contexte de la conversation.} Quand un parcours attend une information (un montant,
une durée), le message est d'abord lu comme une réponse à cette question : \og 250\,000 \fg{} n'est
pas une nouvelle demande. Le classifieur n'est consulté que si le message ne répond pas à la
question en cours.

\paragraph{Les entités, par règles.} Montants (\og 500\eur{} \fg, \og 1\,900 euros \fg, \og 250k
\fg), durées (\og 10 ans \fg, \og 48 mois \fg), liens (\og mon fils \fg), objets (voiture,
appartement, études), dates relatives (\og hier soir \fg) : des expressions régulières suffisent et
coûtent zéro. Si une information manque après les règles, on ne la devine pas avec un LLM : le
parcours la demande au client. Un appel au LLM de moins, et aucun risque d'erreur d'extraction.

\paragraph{L'insatisfaction.} Deux signaux combinés : le classifieur (intention
\f{insatisfaction}) et l'historique (nombre de contacts du jour). Au troisième contact, la règle de
passage de main s'applique même si le ton est neutre.

\section{Étage 2 --- Converser}

\paragraph{Un parcours par intention de conseil ou d'action.} Chaque parcours est un fichier YAML
qui déclare les informations nécessaires, la question à poser pour chacune, l'ordre, et ce qui
se passe à la fin (un conseil, une action, un passage de main). Le gestionnaire de dialogue est un
automate qui lit ces fichiers : déterministe, testable, et modifiable sans toucher au code.

\begin{code}
\footnotesize
\begin{verbatim}
# src/converser/parcours/credit_immobilier.yaml (extrait)
famille: credit
infos:
  revenus_nets:   {source: profil, sinon: "Quels sont vos revenus nets mensuels ?"}
  credits_cours:  {source: profil, sinon: "Avez-vous des credits en cours ?"}
  prix_bien:      {question: "Quel est le prix du bien que vous visez ?"}
  apport:         {source: profil.epargne, confirmer: true,
                   question: "Quelle part de votre epargne souhaitez-vous y consacrer ?"}
fin: {conseil: credit_immobilier, mention_legale: credit, proposer: rdv_conseiller}
\end{verbatim}
\end{code}

\paragraph{Ne jamais demander ce qu'on sait.} C'est la forme la plus visible de la
personnalisation (E4) : si le profil connaît les revenus, l'assistant \emph{confirme} au lieu de
demander (\og J'ai noté vos revenus de 2\,900\eur{} nets par mois \fg). Une information simulée
est confirmée, jamais affirmée.

\paragraph{Le registre.} Un mode \og mots simples \fg{} s'active de lui-même pour un client âgé de
70 ans ou plus ou signalé \og peu à l'aise en ligne \fg{} (Paul), et à la demande pour tout le
monde : phrases courtes, une seule étape par message, aucun sigle sans explication. Ce n'est pas un
autre modèle, c'est un autre jeu de gabarits et une consigne de plus au petit LLM.

\paragraph{Le jargon.} Le glossaire (annexe~\ref{an:glossaire}) est une liste de fiches courtes,
rédigées et vérifiées par nous. Quand un terme du glossaire apparaît dans une réponse, il est
souligné et explicable d'un clic ; quand le client demande \og c'est quoi \dots{} \fg, la fiche est
retrouvée par similarité de vecteurs, en mémoire, sans LLM.

\section{Étage 3 --- Conseiller}

\paragraph{Trois passes.} Le moteur de conseil applique d'abord les conditions
\textbf{d'éligibilité} (dures : âge, revenu, plafond déjà atteint), puis des règles de
\textbf{pertinence} (horizon, objectif, épargne de précaution existante), puis un
\textbf{classement}. Chaque règle a un identifiant et une phrase d'explication paramétrée : quand
elle s'applique, la phrase remplie entre dans les raisons de la proposition.

\paragraph{Pourquoi, et pourquoi pas.} La sortie contient les produits proposés avec leurs raisons
\emph{et} les produits écartés avec la règle qui les écarte. C'est l'explication la plus
convaincante qui soit pour un client --- \og le LEP n'est pas proposé parce que votre revenu fiscal
dépasse le plafond \fg{} --- et pour le jury (15\,\% de la note).

\paragraph{L'appétence pour l'épargne.} Bank Marketing dit, pour 45\,211 clients, qui a souscrit un
dépôt à terme. Une régression logistique entraînée sur les seules caractéristiques du client (âge,
métier, situation, études, solde, crédits, campagnes précédentes ; \textbf{sans} \f{duration})
donne une probabilité d'intérêt pour l'épargne. Elle sert à \emph{ordonner} les propositions
d'épargne et à expliquer cet ordre par les trois facteurs qui pèsent le plus. On mesure l'aire sous
la courbe ROC sur un jeu de test tenu à part ; LightGBM ne remplace la régression logistique que si
le gain est mesuré. Les limites sont dites : données portugaises de 2008 à 2010, donc un signal
faible et illustratif, jamais une décision.

\paragraph{La prudence sur le crédit.} Pour un crédit immobilier, les règles sont celles du Haut
Conseil de stabilité financière : \textbf{mensualités au plus égales à 35\,\% des revenus nets,
assurance comprise, et durée au plus de 25 ans}. La capacité d'emprunt se calcule à partir de la
mensualité maximale $m$, du taux mensuel $i$ et du nombre de mois $n$ :
\[
  m = 0{,}35 \times R - C, \qquad
  \text{capacité} = m \times \frac{1 - (1+i)^{-n}}{i}
\]
où $R$ est le revenu net mensuel et $C$ les mensualités déjà en cours. Le taux est un
\emph{paramètre daté} du catalogue, affiché comme \og taux indicatif \fg. Pour le crédit à la
consommation, la même logique de taux d'effort s'applique, avec la mention légale obligatoire :
\og Un crédit vous engage et doit être remboursé. Vérifiez vos capacités de remboursement avant de
vous engager. \fg{} L'assistant ne promet jamais un accord : il donne une estimation et propose un
rendez-vous avec un conseiller, qui décide.

\section{Étage 4 --- Agir}

L'oral laissait le choix entre guider et agir (E13). Nous agissons, mais avec une règle reprise de
Timso AI : \textbf{la lecture est libre, l'effet de bord demande une confirmation explicite}, et
chaque action laisse une trace dans un journal.

\begin{center}\small
\begin{tabularx}{\textwidth}{@{}l Y c Y@{}}
\toprule
\textbf{Outil} & \textbf{Ce qu'il fait (sur le SI simulé)} & \textbf{Confirmation} &
  \textbf{Rend} \\
\midrule
\f{consulter\_profil} & Lit le profil et l'historique du client & non & Le profil \\
\f{expliquer\_frais} & Retrouve les frais du mois dans l'historique et leur cause & non & Liste
  datée des frais, avec la règle qui les explique \\
\f{bloquer\_carte} & Met la carte en opposition & \textbf{oui} & Une référence d'opposition,
  l'heure \\
\f{preparer\_virement} & Prépare un virement vers un bénéficiaire enregistré ; vérifie le solde ;
  ne l'exécute qu'après accord & \textbf{oui} & Un récapitulatif, puis une référence \\
\f{transmettre\_conseiller} & Crée une demande de rappel avec un résumé de la conversation & oui &
  Un numéro de ticket \\
\bottomrule
\end{tabularx}
\end{center}

Le résumé transmis au conseiller est construit à partir des faits du parcours (intention, infos
recueillies, propositions) : le conseiller n'a pas à reposer les questions. C'est le
\og brouillon \fg{} de Timso AI transposé : l'humain reprend la main avec le travail déjà fait.

\section{Étage 5 --- Formuler}

\paragraph{Le gabarit d'abord.} Pour la plupart des intentions de service (carte, frais, virement,
compte), une réponse à trous, écrite à l'avance et relue, est meilleure qu'une génération : exacte,
instantanée, gratuite. Le petit LLM n'intervient que lorsque la réponse dépend de plusieurs faits à
articuler (un conseil, une explication de refus).

\paragraph{Un seul appel, à la fin.} C'est le seul endroit où un LLM intervient dans le cas
normal : une fois par tour, après que tout a été compris, décidé et calculé. L'escalade de
l'étage 1 est l'exception, comptée et affichée.

\paragraph{Ce que reçoit le petit LLM.} Une consigne système courte (registre de la banque,
vouvoiement, interdits), les \emph{faits structurés} produits par les étages 3 et 4 (propositions,
raisons, chiffres), et la question en cours. Jamais le catalogue entier, jamais l'historique brut :
le contexte reste petit, donc rapide et peu coûteux.

\paragraph{Le garde-fou.} Après génération, trois contrôles automatiques : chaque nombre du texte
doit figurer dans les faits fournis ; aucune formule interdite (\og accordé \fg, \og garanti \fg,
\og sans risque \fg) ; longueur bornée. En cas d'échec, la réponse gabarit est affichée à la place
et l'échec est compté. Le taux d'échec du garde-fou fait partie des résultats présentés.

\paragraph{Le modèle.} \f{ministral-3:3b} par défaut (3,0 Go en Q4\_K\_M) : un modèle européen,
conçu par Mistral AI, à l'aise en français --- ce qui rejoint la question de la souveraineté que
CGI a soulevée. Mais ce choix n'est pas un a priori : le chapitre~\ref{ch:modeles} décrit comment
il est confirmé ou remplacé par la mesure.

\section{Mesurer, en continu}

Chaque appel à \f{reply()} renvoie une trace : le niveau qui a compris la demande (mots-clés,
classifieur ou escalade), puis, pour chaque étage traversé, son nom, le modèle utilisé, la latence, les tokens consommés (Ollama renvoie \f{prompt\_eval\_count} et
\f{eval\_count}) et, pendant les campagnes de mesure, l'énergie. L'interface affiche cette trace en
direct dans le panneau \og frugalité \fg{} : le jury voit, pendant la démo, qu'une opposition de
carte n'a appelé aucun LLM. Le protocole de mesure de l'énergie est au chapitre~\ref{ch:evaluation}.

\section{Un exemple complet : l'appartement}

Léa (C07), 31 ans, 2\,900\eur{} nets par mois, aucun crédit en cours, 25\,000\eur{} d'épargne.

\begin{center}\small
\begin{tabularx}{\textwidth}{@{}l Y l@{}}
\toprule
\textbf{Étape} & \textbf{Ce qui se passe} & \textbf{Modèle} \\
\midrule
\multicolumn{3}{@{}l}{\textit{Tour 1 --- \og Je veux acheter un appartement \fg}} \\
Comprendre & Famille \f{credit}, intention \f{credit\_immobilier}. Les mots-clés \og acheter \fg{}
  et \og appartement \fg{} désignent une seule famille, dans une phrase courte : le filtre conclut
  seul. Entité \f{objet=appartement} & mots-clés \\
Converser & Parcours \f{credit\_immobilier}. Revenus et crédits connus par le profil ; épargne
  connue, à confirmer comme apport. Manque : le prix du bien & --- \\
Formuler & \og C'est un beau projet. J'ai noté vos revenus de 2\,900\eur{} nets par mois. Quel
  est le prix du bien que vous visez, et quelle part de vos 25\,000\eur{} d'épargne souhaitez-vous y
  consacrer ? \fg & gabarit \\
\midrule
\multicolumn{3}{@{}l}{\textit{Tour 2 --- \og Autour de 250\,000, et je peux mettre les 25\,000 \fg}} \\
Comprendre & Lu comme la réponse attendue : prix 250\,000\eur{}, apport 25\,000\eur{} & règles \\
Conseiller & Frais de notaire estimés à 7,5\,\% dans l'ancien : 18\,750\eur{}. Besoin de
  financement : 243\,750\eur{}. À un taux indicatif de 3,5\,\% sur 25 ans, mensualité d'environ
  1\,220\eur{}, soit \textbf{42,1\,\% des revenus : au-dessus des 35\,\%}. Prix maximal finançable
  avec cet apport : environ 212\,000\eur{} & règles \\
Formuler & Explique le calcul, propose trois leviers (bien moins cher, apport plus élevé,
  co-emprunteur), signale qu'un conseiller peut étudier les aides (PTZ selon conditions), propose
  un rendez-vous & \f{ministral-3:3b} \\
Garde-fou & Tous les nombres du texte figurent dans les faits ; aucune promesse & règles \\
\bottomrule
\end{tabularx}
\end{center}

Le panneau \og pourquoi \fg{} affiche alors les règles appliquées (taux d'effort maximal de
35\,\%, durée maximale de 25 ans, estimation des frais de notaire), le taux indicatif avec sa date
et sa source, et la ligne écartée : \og prêt sur 30 ans : au-delà de la durée maximale \fg. Le
calcul est vérifiable à la main par le jury.

\section{Les parcours}

Neuf parcours couvrent les neuf situations du sujet ; les familles de service restantes de
Banking77 reçoivent une fiche de réponse courte et une orientation. Le détail de chaque parcours
est en annexe~\ref{an:parcours}.

\begin{center}\small
\begin{tabularx}{\textwidth}{@{}l Y Y@{}}
\toprule
\textbf{Parcours} & \textbf{Ce que l'assistant demande} & \textbf{Comment il conclut} \\
\midrule
Carte perdue (Inès) & Quelle carte ; perte ou vol & \f{bloquer\_carte} après confirmation ; rassure
  sur la responsabilité ; propose une nouvelle carte \\
Frais (Marc) & Quel frais, si plusieurs & \f{expliquer\_frais} sur l'historique ; plafond légal ;
  si les incidents se répètent, oriente vers l'offre clientèle fragile \\
Virement (Paul) & Bénéficiaire, montant, date & \f{preparer\_virement} pas à pas ; rappel :
  la banque ne demande jamais un code \\
Épargne (Sophie) & Pour qui, horizon, montant mensuel & Propositions classées, raisons, produits
  écartés \\
Appartement (Léa) & Prix, apport & Capacité, taux d'effort, leviers, rendez-vous \\
Crédit auto (Karim) & Montant, durée, apport & Taux d'effort avec le crédit en cours, mention
  légale, rendez-vous \\
Prêt étudiant (Lucas) & Études, montant, caution & Orientation vers le prêt étudiant ; Livret Jeune
  pour épargner \\
Jargon & Le terme & Fiche du glossaire \\
Insatisfaction (Julie) & Rien : on écoute & Reconnaît la frustration, s'excuse,
  \f{transmettre\_conseiller} avec le résumé des trois contacts \\
\bottomrule
\end{tabularx}
\end{center}

\section{Conformité}

\begin{itemize}
  \item \textbf{Données personnelles.} Toutes les données sont simulées ; tout est traité en local,
        aucune donnée ne sort de la machine (E8).
  \item \textbf{Règlement européen sur l'IA.} Les systèmes destinés à évaluer la solvabilité des
        personnes physiques sont classés à haut risque. Notre assistant ne note pas la solvabilité :
        il applique des règles publiques, donne une estimation indicative, et renvoie la décision à
        un conseiller. Le modèle d'appétence n'ordonne que des produits d'épargne. Le règlement
        impose aussi d'informer l'utilisateur qu'il échange avec une IA : un bandeau permanent le
        dit.
  \item \textbf{Crédit.} La mention légale accompagne toute réponse sur un crédit.
  \item \textbf{Protection.} Toute réponse qui touche à la carte ou au virement rappelle qu'aucun
        code ne doit être communiqué --- une protection simple contre l'arnaque au faux conseiller,
        qui vise en particulier les clients âgés.
\end{itemize}
